% !PRINTED 727
\weberpage{727}

%<<CONT
formula (9) reduce all the values of $\beta$ to the interval from $0$ to
$\nu' - 1$, then we obtain all the numbers of this interval.

Accordingly formula (3) can be represented thus:
\weq{10}{
X_1 = -\frac{2}{m}\,(1,\,\Theta^{-1},\,r)\sum_{0,\,\nu'-1}^{\beta}\Theta^\beta\lambda_\beta.
}
This we have to substitute in the expression given in \S.\,197, (8)
\[
D B = 2^{2^{\varkappa-4}(\varkappa-2)}\prod X_1 = m^{\nu''}2^{-\nu'}\prod X_1
\]
and then, if we set
\weq{11}{
U_\Theta = \sum_{0,\,\nu'-1}^{\beta}\Theta^\beta\lambda_\beta
}
there results, with regard to \S.\,197, (14):
\weq{12}{
D B = \prod^{\Theta} U_\Theta.
}
This product can now be represented in the following way by a
determinant which contains only the conjugate logarithms of the
units $\tau$.

If we omit the index $\Theta$ on $U$, we can set up the system of
equations:
\[
\begin{array}{l}
U \;=\; \lambda_0 \;+\; \lambda_1\Theta \;+\; \lambda_2\Theta^2 \;+\; \cdots \;+\; \lambda_{\nu'-1}\Theta^{\nu'-1} \\[4pt]
\Theta U \;=\; -\lambda_{\nu'-1} \;+\; \lambda_0\Theta \;+\; \lambda_1\Theta^2 \;+\; \cdots \;+\; \lambda_{\nu'-2}\Theta^{\nu'-1} \\[4pt]
\text{\normalfont .\;.\;.\;.\;.\;.\;.\;.\;.\;.\;.\;.\;.\;.\;.\;.\;.\;.\;.\;.\;.\;.\;.\;.\;.\;.\;.\;.\;.\;.} \\[4pt]
\Theta^{\nu'-1} U \;=\; -\lambda_1 \;-\; \lambda_2\Theta \;-\; \lambda_3\Theta^2 \;+\; \cdots \;+\; \lambda_0\Theta^{\nu'-1},
\end{array}
\]
from which, by elimination of the powers of $\Theta$, it follows that:
\weq{13}{
\begin{vmatrix}
\lambda_0 - U, & \lambda_1, & \lambda_2, & \ldots, & \lambda_{\nu'-1} \\
-\lambda_{\nu'-1}, & \lambda_0 - U, & \lambda_1, & \ldots, & \lambda_{\nu'-2} \\
\multicolumn{5}{c}{\text{\normalfont .\;.\;.\;.\;.\;.\;.\;.\;.\;.\;.\;.\;.\;.\;.\;.}} \\
-\lambda_1, & -\lambda_2, & -\lambda_3, & \ldots, & \lambda_0 - U
\end{vmatrix} = 0.
}
% [source note: the first row prints a semicolon-shaped mark after \lambda_0-U; the mathematically required list punctuation is normalized to a comma.]
This is, for $U$, an equation of degree $\nu'$, whose roots are the
quantities $U_\Theta$. If we arrange it according to powers of $U$,
then, since $\nu'$ is even, the highest power $U^{\nu'}$ of $U$ receives
the coefficient $1$, and one therefore finds the product of the roots
occurring in (12) from the determinant (13) by setting $U = 0$.

The determinant so obtained we now wish to arrange somewhat
differently, by leaving the first row in its place, while the
%CONT>>
